\section{2024-10-25 - de-Broglie - Wellenlänge}

\startitemize[packed]
\item 1924: aus theoretischen Überlegungen
\stopitemize

\startformula
\lambda = \frac{h}{p} \leftrightarrow p = \frac{h}{\lambda}
\stopformula

mit $\lambda$ : Wellenlänge und $p$ : Impuls

Dabei gilt $p=m \cdot v$ mit (Masse $m$, Geschwindigkeit $v$).
$h$ ist de Planck'sche Konstante mit:

\startformula
h = 6,626 \cdot 10^{-34} Js
\stopformula

Formal gilt:

\startformula
\frac{dE}{dv}=p \qquad (\text{E Energie})
\stopformula

Damit ergibt sich:

\startformula
\lambda(v)=\frac{h}{m\cdot v}=\frac{h}{m}\cdot v^{-1}
\stopformula

\startframedtask
Vergleichen Sie den Vorfaktor der Regression $0,0067\frac{m^2}{s}$ mit dem theoretisch erwarteten Wert.
\stopframedtask

\startformula
\frac{h}{m}\approx 7,27 \cdot 10^{-4}
\stopformula

\startformula
\[\frac{h}{m}\] = \frac{Js}{Kg}=\frac{kg \cdot \frac{m^2}{s^2}\cdot s}{kg}=\frac{kg \cdot \frac{m^2}{s}}{kg}=\frac{m^2}{s}
\stopformula

Unsere Überprüfungs...

hat als Mittelwert für $\lambda \cdot v \approx 7,66 \cdot 10^{-4}\frac{m^2}{s}$ bzw. $\lambda \cdot v \approx = \frac{h}{m}$ egeben:

% \startformula
% \lambda = \frac{h}{m\cdot v} \LeftRightarrow \lambda \cdot v = \frac{h}{m}
% \stopformula

Gemäß der Ringe 1 ergibt sich:

% \startformula
% \lambda = 6,66\cdot 10^{-4}\frac{m^2}{s} \cdot \frac{1}{v}
% \stopformula


Die Messung passt also gut

\blank
\hrule
\blank

\startframedtipp
{\bf Spektrum schwarzer Strahler}
\stopframedtipp

\startframedtipp

Puls ist ableitung der Energie

\stopframedtipp

\startframedtipp
Elektronenkanone ebenfalls erläutern
\stopframedtipp

